\documentstyle[11pt]{article}
\setlength{\textwidth}{170mm}
\setlength{\textheight}{220mm}
\setlength{\oddsidemargin}{0mm}
\setlength{\evensidemargin}{10mm}
\setlength{\topmargin}{-10mm}
\def\RR{\hbox{\raise.4ex\hbox{$\scriptstyle{|}$}\hskip-1.85pt R}}
%\renewcommand{\baselinestretch}{1.2}
\newcommand{\REF}{\nopagebreak\noindent{\large{\bf
      References}}\vspace{.2in}}

\newcommand{\Ref}[1]{\medskip\noindent\hangafter=1
      \hangindent=20pt\ignorespaces #1}
\pagestyle{empty}
\begin{document}


\noindent
University of Illinois \hfill Department of Economics \\
Spring 1998 \hfill Econ 476 \hfill Roger Koenker

%\vspace{4mm}
\begin{center}
\large{
{\bf Advanced Econometrics I \\
\vspace{2mm}
Course Prospectus}}
\end{center}

\vspace{10mm}
The field sequence 476-477 in econometrics is designed to give students a 
working knowledge of a broad array of current topics in econometric theory and 
prepare them for empirical applications near the research frontier.  In the 
first course, 476, we intend to focus the first half of the course on 
asymptotic methods and the study of the large-sample behavior of econometric 
procedures supplemented by an introduction to monte carlo 
methods.  The second half of the course will be devoted to recent developments 
in semi-parametric methods in econometrics.

\vspace{5mm}
\noindent
{\bf Suggested Text:}

Gourieroux, C. and A. Monfort (1995) 
{\it Statistics and Econometric Models}, Cambridge U. Press.

\vspace{5mm}
\noindent
{\bf Recommended Supplementary Tests:}

Davidson, R. and MacKinnon J.G. (1993), {\em Estimation and Inference in 
Econometrics}, Oxford U. Press.

McCabe, B. and Tremayne, A. (1993), {\em Elements of Modern Asymptotic 
Theory with Statistical Applications}, Manchester U. Press.

\vspace{5mm}
\noindent
Readings for the second half of the course will be distributed in class.
Regular problem sets, a brief research paper and final exam will be 
required.  The problem sets will attempt to intermingle theoretical, 
empirical and monte-carlo exercises.

\vspace{10mm}
\begin{center}
{\bf Outline}
\end{center}

\begin{enumerate}
\item {\bf Introduction to Asymptotic Theory and Monte-Carlo Methods}

Asymptotic theory may be  the violin of econometric analysis and 
monte-carlo  the timpani. But this is not  to disparage either;  both are 
indispensable instruments required to explore the behavior of the 
wide array of estimation and inference methods which are currently 
employed in econometrics.  
Equally indispensible are the remaining instruments of orchestra with
which we may identify the richness of empirical applications in economics.
Asymptotics and monte carlo are complementary.  To the 
extent that asymptotic theory is a reliable guide to the evaluation of the 
finite sample performance of estimators and test statistics,  This can 
only be confirmed by appropriately designed monte-carlo experiments.  Thus,
asymptotics provides guidance on the design of simulation experiments, 
while the experiments serve to evaluate the adequacy of the approximations 
suggested by the asymptotics.

In the first part of the course we will attempt to interweave 
asymptotics and monte-carlo introducing the basic techniques of both 
approaches and illustrating their application on a wide variety of problems
of estimation and inference including

\begin{enumerate}
\item (generalized) maximum likelihood and related methods of inference 
 \item location/scale estimation 
 \item (generalized) linear models (exponential family models) 
 \item (quasi) maximum likelihood and related methods of inference
\end{enumerate}

\noindent
Along the way we will revisit, or encounter for the first time in some 
cases:  modes of convergence, laws of large numbers, central 
limit theorems, stochastic equicontinuity, and weak convergence.  The basic
tools of the monte-carlo trade will also be introduced including some attention
to design of experiments, variance reduction techniques and a brief
introduction to density estimation.  Particular attention will be paid to
monte carlo methods for evaluation of inference procedures.  

\item {\bf Semi-Parametric Methods in Econometrics}

The second half of the course will be devoted to recent developments in 
semi-parametric methods of estimation and inference.  These methods seek to 
relax the stringent assumptions required to justify the strictly parametric 
models employed in classical econometrics.  We will consider a selection of
the following topics.

\begin{enumerate}
\item Nonparametric density estimation
\item Nonparametric regression
\item Quantile regression
\item Semiparametric discrete choice models
\item The bootstrap and related resampling methods
\item Measures of inequality and tail behavior
\item Semiparametric survival analysis
\item Transformation models and ``average derivative" estimation
\end{enumerate}


\end{enumerate}
\end{document}
